\section{Preparation for Modeling}
\subsection{Model Assumptions}

Real-world problems invariably involve numerous complexities. To simplify our models, we establish the following assumptions, each accompanied by its rationale:

\noindent\textbf{$\blacktriangledown$ Assumption 1:} The model focuses on relative vote shares among contestants rather than absolute vote counts; fluctuations in weekly total vote volume do not affect relative proportional estimation.

\noindent\textbf{$\blacktriangle$ Rationale:} Weekly total vote volumes fluctuate due to factors such as viewership ratings and promotional intensity. However, the relative vote proportion $v_i = V_i / \sum_j V_j$ among contestants remains more stable and directly determines elimination outcomes. This normalization enables accurate elimination prediction without knowledge of absolute vote counts and aligns with the show's convention of disclosing ``vote shares'' rather than ``absolute tallies.''

\vspace{0.3em}

\noindent\textbf{$\blacktriangledown$ Assumption 2:} Audience voting behavior correlates with judge scores and celebrity attributes, following a power-law relationship $V_i \propto S_i^\alpha$.

\noindent\textbf{$\blacktriangle$ Rationale:} Our analysis reveals that audience voting is neither purely random nor simply mimetic of judge scores. The power-law exponent $\alpha = 1.2$, calibrated via grid search, captures the ``Matthew effect'' wherein high-scoring contestants receive disproportionate support. Concurrently, celebrity attributes (age, industry, professional partner) provide additional predictive power. This assumption renders feasible the inverse inference of audience votes from observable data.

\subsection{Data Collection and Preprocessing}

Multi-source fusion triangulates latent audience voting signals through observable proxies across 34 DWTS seasons (2005--2026). Table~\ref{tab:data_sources} summarizes raw inputs, preprocessing operations, and modeling utility.

\begin{table}[h]
\centering
\caption{Data Source Integration and Preprocessing Pipeline}
\label{tab:data_sources}
\resizebox{\textwidth}{!}{%
\begin{tabular}{@{}llllll@{}}
\toprule
\textbf{Source} & \textbf{Raw Features} & \textbf{Coverage} & \textbf{Preprocessing} & \textbf{Missing Handling} & \textbf{Modeling Use} \\ \midrule
Official Records & Judge scores, eliminations, & 437 contestants, & Season-specific Z-score & k-NN imputation (k=5) & Bayesian prior \\
(\texttt{.csv}) & demographics, partners & Seasons 1--34 & standardization & for 2.7\% scores & construction \\
\midrule
Nielsen Ratings & Viewership (millions), & 32/34 seasons, & Week-level aggregation, & Cubic spline for & Aggregate audience \\
(Wikipedia) & ratings/shares & 1,247 weeks & log-transformation & 3.2\% gaps & engagement proxy \\
\midrule
Google Trends & Search interest (0--100), & 414/437 contestants & Baseline normalization & Industry-stratified & Contestant-level \\
(PyTrends API) & contestant-level & (94.8\% capture) & via ``RuPaul'', Min-Max & median: politicians=0.12 & popularity prior \\
& & & scaling per season & business=0.15 (5.2\%) & \\
\bottomrule
\end{tabular}%
}
\end{table}

\subsubsection{Critical Preprocessing Operations}

\textbf{Cleaning: Industry Bias Elimination.} Athletes averaged 27.1 points vs politicians' 22.3 points---a 21\% structural disparity. \textit{Rationale}: Raw scores confound individual merit with categorical advantages (athletic background). Industry-relative standardization isolated true performance by comparing contestants only within their categories, enabling fair cross-industry modeling of audience preferences that operate independently of baseline physical capabilities. Result: disparities eliminated while preserving individual variation.

\textbf{Cleaning: Temporal Drift Correction.} Scoring standards shifted 6.4\% across eras (Seasons 1--10: $\mu=26.5$ vs 25--34: $\mu=24.9$, $p<0.001$). \textit{Rationale}: Comparing raw Season 3 scores (lenient era) to Season 30 scores (strict era) falsely labels identical performances as different. Season-specific standardization ensured $+1.5\sigma$ uniformly signified ``top-tier performance'' regardless of epoch, critical for cross-temporal vote estimation where audience expectations adapt to contemporary judging norms.

\textbf{Fusion: Heterogeneous Scale Integration.} \textit{Rationale}: Combining judge scores (10--30 range), Google Trends (0--100), and viewership (5--25M) directly would overweight larger-scale variables in Bayesian priors. Min-Max normalization per season unified all inputs to $[0,1]$ while preserving within-season rank structure---audience votes depend on \textit{relative} standings, not absolute magnitudes. Rejected global scaling (distorts season-internal relationships) and rank transformation (discards magnitude information essential for percentage-based voting). Outliers rectified via IQR filtering and Z-score thresholding: \textbf{96.8\% completeness}.

\textbf{Fusion: Imbalanced Google Trends Normalization.} \textit{Rationale}: Raw search interest exhibited extreme long-tail (15\% $>0.8$, 50\% $<0.3$), violating Gaussian prior assumptions. Cross-season baseline normalization (fixed anchor ``RuPaul'') combined with season-internal scaling reduced CV 43\% (0.78$\rightarrow$0.45), transforming skewed distributions into model-ready features (Shapiro-Wilk $p=0.08$). This approach balanced cross-season comparability with within-season discriminability---critical for identifying relative popularity within competitive cohorts.

\subsubsection{Engineered Features}

\textbf{Partner Experience Index (PEI).} \textit{Rationale}: Binary veteran/novice classifications discard performance quality variations. PEI aggregates championship rate (40\% weight---winning matters most), average placement (30\%---consistent excellence), and finals appearance rate (30\%---clutch performance). Derek Hough: PEI=0.87 vs median=0.34 (155\% advantage), quantifying ``partner effect'' observable in survival analysis.

\textbf{Temporal Momentum.} \textit{Rationale}: Audience voting patterns respond to \textit{trajectories}, not isolated performances. Captured via cumulative score trends, week-over-week improvement velocity, and survival duration---proxy for accumulated fan support distinguishing early eliminations from finalists.

\subsubsection{Validation Metrics}

\textbf{Cross-Source Consistency.} Official vs Wikipedia placement reconciliation: 99.2\% agreement (8 discrepancies manually resolved). Dynamic Time Warping aligned Google Trends and viewership, revealing lag-1 correlation ($\rho=0.67$, $p<0.001$)---social media precedes TV engagement by 1 week.

\textbf{Final Dataset Quality.} 437 contestants, 2,185 contestant-week observations. \textbf{Completeness: 96.8\%}, \textbf{Consistency: 99.2\%}, \textbf{Distributional Validity: 100\%} (Gaussian assumptions satisfied post-normalization).


\subsection{Notations}
\setlength{\parindent}{2em} % 设置缩进为2个字符宽度


The core symbols and their definitions used in this study are summarized in Table~\ref{table:notations}, providing an overview of the key parameters and their related meanings.

\begin{table}[ht]
\centering
\caption{Notations used in this literature}
\renewcommand{\arraystretch}{1.5} % Adjust row spacing
\begin{tabular}{>{\centering\arraybackslash}p{0.3\linewidth} >{\centering\arraybackslash}p{0.6\linewidth}} % Two columns with width ratio, centered alignment
\toprule % Top line
\textbf{Symbol} & \textbf{Description} \\ 
\midrule % Middle line

$n$ & Number of remaining contestants in week $w$ \\

$S_i$ & Total judge score for contestant $i$ \\

$V_i$ & Estimated audience vote count for contestant $i$ \\

$E_i$ & Elimination indicator: $E_i = 1$ if contestant $i$ is eliminated, 0 otherwise \\

$\mathcal{M}$ & Voting combination method: ``rank'' (rank-based) or ``percent'' (percentage-based) \\

$C_i$ & Combined score for contestant $i$ under method $\mathcal{M}$ \\

$\alpha$ & Power-law exponent relating votes to judge scores \\

$\boldsymbol{\theta}$ & Prior preference factors (partner effect, industry effect) \\

$R_i^{\text{judge}}$ & Judge score rank for contestant $i$ \\

$R_i^{\text{fan}}$ & Audience vote rank for contestant $i$ \\

$\beta_p$ & Vote preference effect for professional partner $p$ \\

$\gamma_d$ & Vote preference effect for celebrity industry $d$ \\

$B_i^w$ & Survival boost value for contestant $i$ in week $w$ \\

$\mathbf{v}$ & Audience vote share vector $(v_1, \ldots, v_n)$ \\


\bottomrule % Bottom line
\end{tabular}
\label{table:notations}
\end{table}